\section{Appendix}

\subsection{Patent benchmark specification}

Because the main paper depends on \texttt{BindingDB\_patent / patent\_temporal}, we make its construction explicit here. Table~\ref{tab:patent-spec} enumerates the year boundaries, split counts, label construction, overlap policy, and seed semantics for both the main patent split and the earlier-cutoff robustness variant.

\begin{table}[t]
  \centering
  \scriptsize
  \caption{Explicit specification of the primary patent temporal/provenance benchmarks. The temporal partition is fixed by year; the random seed changes model stochasticity, not the partition itself.}
  \label{tab:patent-spec}
  \setlength{\tabcolsep}{3.5pt}
  \resizebox{\linewidth}{!}{
  \begin{tabular}{@{}lll@{}}
    \toprule
    Aspect & \texttt{patent\_temporal} & \texttt{patent\_temporal\_v2017} \\
    \midrule
    Source files & \texttt{train\_val.csv + test.csv} & \texttt{train\_val.csv + test.csv} \\
    Train years & 2013--2017 & 2013--2016 \\
    Validation year & 2018 & 2017 \\
    Test years & 2019--2021 & 2018--2021 \\
    Rows (train / valid / test) & 142,275 / 41,155 / 49,028 & 90,771 / 51,504 / 90,183 \\
    Pos. rate (train / valid / test) & 0.566 / 0.534 / 0.575 & 0.568 / 0.562 / 0.556 \\
    Label transform & \multicolumn{2}{l}{raw patent value $Y \mapsto \mathrm{p}K_d = 9 - Y / \ln(10)$; positive label if $\mathrm{p}K_d \ge 7.0$} \\
    Overlap policy & \multicolumn{2}{l}{no cold-entity filtering; entity overlap is allowed, exported, and diagnosed via \texttt{overlap\_score} buckets} \\
    Duplicate policy & \multicolumn{2}{l}{pair rows are kept after year partition; entity export tables are deduplicated by canonical ID} \\
    Seed semantics & \multicolumn{2}{l}{partition is fixed across seeds; seed affects training / initialization only} \\
    OOD interpretation & \multicolumn{2}{l}{temporal/provenance OOD rather than pure time-only shift} \\
    \bottomrule
  \end{tabular}}
\end{table}


\subsection{Additional real-OOD support benchmarks}

The main paper focuses on the strongest evidence chain: the synthetic \texttt{BindingDB\_Kd} reference panel, the 3-seed \texttt{BindingDB\_patent} ranking reversal, and the patent year/overlap diagnosis. \texttt{BindingDB\_Ki} and \texttt{DAVIS} remain supporting real-OOD benchmarks: both now use matched 3-seed evaluation, but they reinforce rather than replace the patent temporal story.

\begin{table}[t]
  \centering
  \scriptsize
  \caption{Secondary real-OOD support benchmarks. These panels are seed0-only and therefore serve as supporting counterexamples rather than headline evidence.}
  \label{tab:support}
  \setlength{\tabcolsep}{4pt}
  \begin{tabular}{@{}lccc@{}}
    \toprule
    Benchmark & base (\AUPRC{}/\AUROC{}) & FTM sparse (\AUPRC{}/\AUROC{}) & $\Delta$ vs.\ base \\
    \midrule
    \texttt{BindingDB\_Ki / unseen\_target} &
    0.6574 / 0.7103 &
    0.6295 / 0.6855 &
    -0.0279 / -0.0248 \\
    \texttt{DAVIS / unseen\_target} &
    0.4258 / 0.8849 &
    0.3705 / 0.8757 &
    -0.0554 / -0.0092 \\
    \bottomrule
  \end{tabular}
\end{table}


On \texttt{BindingDB\_Ki / unseen\_target}, \texttt{base} remains the strongest internal model at $0.6934 \pm 0.0257 / 0.7192 \pm 0.0068$, ahead of both \texttt{RAICD} and \texttt{FTM sparse}. On \texttt{DAVIS / unseen\_target}, the same ordering remains: \texttt{base} reaches $0.4097 \pm 0.0275 / 0.8613 \pm 0.0153$, while \texttt{RAICD} and \texttt{FTM sparse} are both lower. Seed0 support screens for the strongest recent baseline \texttt{DTI-LM} are also negative on both datasets, so the synthetic \texttt{unseen\_target} gain does not reliably transfer even after expanding the model family.

\subsection{Paired uncertainty support}

Table~\ref{tab:bootstrap-summary} summarizes the paired bootstrap intervals for every comparison used in a headline ranking claim. The main synthetic gain (\texttt{RAICD} on \texttt{blind\_start}), the synthetic target-cold gain (\texttt{FTM} on \texttt{unseen\_target}), and the real-OOD losses on \texttt{BindingDB\_Ki} and \texttt{DAVIS} all keep their sign under paired resampling. This makes the paper's wording more disciplined: the main and support claims are no longer based only on seed-mean ordering.

\begin{table}[t]
  \centering
  \scriptsize
  \caption{Paired bootstrap support for the comparisons used in headline ranking claims. Positive deltas favor the second model in the comparison; negative deltas favor \texttt{base}.}
  \label{tab:bootstrap-summary}
  \setlength{\tabcolsep}{3.5pt}
  \resizebox{\linewidth}{!}{
  \begin{tabular}{@{}lll@{}}
    \toprule
    Comparison & Mean \AUPRC{} delta & 95\% CI \\
    \midrule
    \texttt{blind\_start}: \texttt{RAICD - base} & $+0.0288$ & [$+0.0145$, $+0.0431$] \\
    \texttt{unseen\_target}: \texttt{FTM - base} & $+0.0214$ & [$+0.0143$, $+0.0286$] \\
    \texttt{patent}: \texttt{RAICD - base} & $-0.0080$ & [$-0.0101$, $-0.0059$] \\
    \texttt{patent}: \texttt{FTM - base} & $-0.0050$ & [$-0.0072$, $-0.0029$] \\
    \texttt{patent}: \texttt{DTI-LM - base} & $+0.0054$ & [$+0.0028$, $+0.0079$] \\
    \texttt{Ki}: \texttt{RAICD - base} & $-0.0249$ & [$-0.0271$, $-0.0226$] \\
    \texttt{Ki}: \texttt{FTM - base} & $-0.0174$ & [$-0.0197$, $-0.0150$] \\
    \texttt{DAVIS}: \texttt{RAICD - base} & $-0.0259$ & [$-0.0375$, $-0.0127$] \\
    \texttt{DAVIS}: \texttt{FTM - base} & $-0.0310$ & [$-0.0434$, $-0.0177$] \\
    \bottomrule
  \end{tabular}}
\end{table}


\subsection{Recent baseline screening and promotion}

To avoid cherry-picking concerns, Table~\ref{tab:recent-screening} summarizes the seed0 screening pass used to decide which recent baselines to promote. The promotion criterion was simple: prioritize the strongest recent model on the primary real-OOD benchmark \texttt{patent\_temporal}, then keep any additional recent baseline that is competitive enough to test whether the main conclusion changes.

\begin{table}[t]
  \centering
  \scriptsize
  \caption{Transparent seed0 screening used to decide which additional baselines were promoted to matched three-seed evaluation. The early external-baseline screen predated the non-patent temporal probe; the later panel-expansion screen added \texttt{RF} and \texttt{GraphDTA-style} under the final five-model design.}
  \label{tab:recent-screening}
  \setlength{\tabcolsep}{3.0pt}
  \resizebox{\linewidth}{!}{
  \begin{tabular}{@{}lccccccp{5.3cm}@{}}
    \toprule
    Model & \texttt{patent\_temporal} & \texttt{nonpatent\_temporal} & \texttt{unseen\_drug} & \texttt{blind\_start} & \texttt{unseen\_target} & Promotion & Rationale \\
    \midrule
    \texttt{RF} &
    0.8249 / 0.7793 &
    0.7216 / 0.7648 &
    0.7279 / 0.8912 &
    0.4488 / 0.7589 &
    0.5618 / 0.7981 &
    main 3-seed promotion &
    strongest seed0 system on every screened temporal and synthetic axis except \texttt{unseen\_target}; therefore promoted across all main axes as the learner-control baseline \\
    \texttt{GraphDTA-style} &
    0.7741 / 0.7242 &
    0.6673 / 0.7444 &
    0.5386 / 0.8159 &
    0.4245 / 0.7426 &
    0.5289 / 0.7776 &
    targeted 3-seed promotion &
    architecture-diverse graph-family probe that is competitive enough on the main synthetic and temporal axes to test whether the expanded-panel conclusion depends on fixed-vector models \\
    \texttt{DTI-LM} &
    0.7897 / 0.7435 &
    -- &
    -- &
    0.3873 / 0.7130 &
    0.5308 / 0.8011 &
    main 3-seed promotion &
    strongest seed0 recent baseline on the primary real-OOD benchmark; therefore promoted in the main text \\
    \texttt{HyperPCM} &
    0.7686 / 0.7247 &
    -- &
    -- &
    0.3896 / 0.7447 &
    0.5294 / 0.7990 &
    secondary 3-seed promotion &
    competitive on synthetic seed0 axes, but weaker on the main patent benchmark; promoted only to test whether the main conclusion changes \\
    \bottomrule
  \end{tabular}}
\end{table}


\subsection{Reusable benchmark resources}

The benchmark release is intended to be reused independently of the in-repo models. Table~\ref{tab:resource-release} lists the concrete artifacts exported for each split and how they interact with the standalone evaluator. In practice, an external user only needs the canonical split bundle plus a prediction CSV keyed by \texttt{example\_id}; the evaluator then reproduces the same global metrics and overlap-bucket metrics reported in this paper.

\begin{table}[t]
  \centering
  \scriptsize
  \caption{Reusable benchmark artifacts released with each split. The release is designed so that external predictions can be scored without re-running the in-repo training code.}
  \label{tab:resource-release}
  \setlength{\tabcolsep}{3.5pt}
  \resizebox{\linewidth}{!}{
  \begin{tabular}{@{}lll@{}}
    \toprule
    Artifact & Example files & Purpose \\
    \midrule
    Split manifest & \texttt{manifest.json} & versioned metadata, split counts, prevalence, seed, and canonical paths \\
    Partition tables & \texttt{train.csv}, \texttt{valid.csv}, \texttt{test.csv} & canonical labeled partitions with \texttt{example\_id} and exported overlap scores \\
    Pair tables & \texttt{train\_pairs.csv}, \texttt{test\_pairs.csv} & lightweight pair-only views for external model pipelines \\
    Entity tables & \texttt{drugs.csv}, \texttt{targets.csv} & unique drug / target vocabularies aligned to the exported split IDs \\
    Full joined table & \texttt{all.csv}, \texttt{all\_pairs.csv} & complete resource snapshot for auditing or custom slicing \\
    Standalone evaluator & \texttt{scripts/evaluate\_prediction\_csv.py} & scores a prediction CSV against the canonical reference split and returns global + overlap-bucket metrics \\
    \bottomrule
  \end{tabular}}
\end{table}


\subsection{Method-centric ablations}

Earlier project stages explored routers, conservative target-memory variants, and selective fallback rules. These analyses were useful for method diagnosis, but they do not define the final paper narrative. The benchmark paper is intentionally centered on ranking instability across shift definitions rather than on proposing a new best-performing method.
